% ============================================================================
% Chapter 4: PokeAgent Baseline Harness
% ============================================================================

\newpage
\section{PokeAgent Baseline Harness}\label{sec:pokeagent}

This chapter traces the development of the \pokeagent{} harness from its simplest form during the \pokeagent{} Challenge NeurIPS 2025 competition through to the realized expert system we deploy as a performance ceiling in the ensuing AutoEvolve experiments in Chapter~\ref{sec:autoevolve}. The narrative follows a chronological arc across three distinct phases of development: (1)~the \emph{competition era} harness--- aptly referred to as \textbf{\pokeagent{} Simplest}---in which a minimal four-module agent participated in the NeurIPS 2025 challenge; (2)~in transitioning from competition to an \emph{open benchmark} we built upon learnings from the simple agent to an intermediate harness which introduced structured objectives, a knowledge base, and A* pathfinding to produce the baseline results reported in the \pokeagent{} Challenge Retrospective~\citep{karten2025pokeagent}; and (3)~the \emph{expert harness} $\mathcal{H}_{\text{expert}}$, a fully realised multi-agent orchestration system engineered for state-of-the-art performance in the long-horizon setting and which provides the human-crafted ceiling against which AutoEvolve is measured in Chapter~\ref{sec:autoevolve}.


% ---------------------------------------------------------------------------
\subsection{Early Development: \pokeagent{} Simple Era}\label{sec:early_development}

The earliest \pokeagent{} harness, deployed as a baseline for the NeurIPS 2025 competition, was a minimal system built to validate whether frontier VLMs could make meaningful progress in Pok\'{e}mon Emerald at all. This system received visual frames and structured state text in accordance with our environment design though it employed no external tools as the agent could only output raw GBA button presses; a rolling short-term buffer of its action history was its only memory abstraction. Finally, it followed an explicitly constructed observe-plan-act loop embedded within its four-module design.

\subsubsection{Four-Module Pipeline}\label{sec:four_module}

\pokeagent{} Simplest implemented a linear pipeline of four modules, three of which backed by a separate VLM call:

\begin{enumerate}[nosep]
    \item \textbf{Perception}: receives the $240 \times 160$ game frame and structured state text; produces a textual scene analysis classifying the current context (overworld, battle, dialog, menu) and a flag indicating whether the situation requires deliberate reasoning.
    \item \textbf{Planning}: receives the memory context and current plan from the previous agent; checks whether the existing plan has been accomplished and, if so, generates an updated plan covering immediate goals, short-term objectives, and long-term strategy.
    \item \textbf{Memory}: involves no external calls to an LLM; maintained purely as a rolling buffer of 50 entries of observation--action pairs. \texttt{deque}. It programatically extracts significant events (map transitions, battle entries/exits) and formats the recent history into a context string, in accordance with pre-defined textual-parsing heuristics.
    \item \textbf{Action}: receives all accumulated context (memory, plan, perception, visual frame, state text) and produces a list of up to 10 Game Boy Advance (GBA) button presses (\texttt{A}, \texttt{B}, \texttt{UP}, \texttt{DOWN}, \texttt{LEFT}, \texttt{RIGHT}, \texttt{START}, \texttt{SELECT}).
\end{enumerate}

Each step through this pipeline required 3 VLM calls, making the agent slow on a per-action basis. Importantly, it introduced a rudimentary objective system at the point of the planning module as it prompted the agent to maintain dynamic objectives that it could create or complete via inline commands (\texttt{ADD\_OBJECTIVE:}, \texttt{COMPLETE\_OBJECTIVE:}). The system prompt was minimal by design with just two sentences describing the agent's role, and no access to MCP tools, long-term memory abstractions in the form of a knowledge base, and no mechanism for self-reflection. The agent navigated purely by interpreting the ASCII map, relying on the foundation model's world knowledge and spatial reasoning to infer appropriate movement. Despite these limitations, it was sufficient to complete the tutorial setting, establishing a proof of concept that frontier VLMs could play Pok\'{e}mon Emerald with minimal scaffolding.


\begin{figure}[t]
\centering
\includegraphics[width=0.85\linewidth]{ch4/simple_agent_draft.png}
\caption{\textbf{PokeAgent Simple: Competition-Era Architecture.} The four-module pipeline (Perception $\to$ Planning $\to$ Memory $\to$ Action) processes a game frame and structured state through sequential VLM calls, producing button presses. The SimpleAgent variant collapses this into a single structured VLM call. Neither variant has access to external tools or persistent memory.}

\tersoo{TODO: Refine this draft figure with a cleaner diagram.}
\label{fig:ch4_simple_agent}
\end{figure}

% ---------------------------------------------------------------------------
\subsection{Early Performance}\label{sec:early_performance}

A fundamental limitation of the four-module pipeline was its perception bottleneck. The Perception module translated a raw $240 \times 160$ game frame into a textual scene description, and all downstream modules (Planning, Memory, Action) operated exclusively on this text. The Action module never saw the original image. This architecture is analogous to a human who must act on an environment they cannot see, relying entirely on a second person's verbal description of the scene---a description that is inevitably lossy, filtered through that person's judgment about what information is relevant. In our setting, the Perception module would routinely omit tile-level details (e.g., ledge positions, NPC facing directions) that proved critical for navigation, while over-emphasizing salient but less actionable features such as the names of visible Pok\'{e}mon species. The result was an information bottleneck that constrained action quality regardless of the Action module's reasoning capability. Later iterations of the harness addressed this by providing the visual frame directly to the action-generating VLM call, eliminating the intermediate textual translation.

\subsubsection{NeurIPS 2025 Competition Results}\label{sec:competition_results}
\begin{figure}[h]
\centering
\includegraphics[width=\linewidth]{ch3/track2_progress_time.pdf}
\caption{\textbf{NeurIPS Speedrunning Leaderboard.} Milestone progress vs.\ wall-clock time (left) and cumulative agent steps (right) for the six teams that completed all 15 milestones. Team Deepest (Judge's Choice) completed in the fewest steps (649) despite ranking 5th by time, highlighting the tradeoff between inference speed and action efficiency. \pokeagent{} Simplest was both time and step-poor.}
\label{fig:ch4_neurips_leaderboard}
\end{figure}

The NeurIPS speedrunning competition (Figure~\ref{fig:ch4_neurips_leaderboard}) produced six teams that achieved 100\% milestone completion, with completion times ranging from 40:13 (Heatz) to 2:46:48. The diversity of approaches is instructive. At the top, Heatz's Scripted Policy Distillation (SPD) method leveraged an LLM to decompose the task into subgoals with executable success conditions, generate scripted policies for each, and distill these into neural networks refined via RL---achieving emergent behaviors including BFS-based local pathfinding and trainer battle skipping~\citep{karten2025pokeagent}. Hamburg Pok\'{e}Runners took second with a pure RL approach: recurrent PPO with a CNN-MLP encoder, LSTM for temporal reasoning, and a binary milestone vector that served as both memory and goal conditioning. Anthonys (third place) combined deterministic A* pathfinding with phase-based prompt templates, decomposing the game into seven distinct phases with specialized navigation strategies. Deepest received the Judge's Choice award for sample efficiency, completing the route in just 649 actions (fewest of any team) using a training-free architecture built around a guidebook system, image-based goal recognition, and adaptive reasoning budgets. Full methodology descriptions for each team are provided in Appendix~\ref{sec:appendix_competitor_methods}.


A key insight from the speedrunning competition is the consistency that specialist methods display in outperforming generalist LLM approaches in the tutorial setting. The winning team, Heatz, completed the route in 40:13---nearly twice as fast as the second-place Hamburg Pok\'{e}Runners, who also employed RL-based approach. The remaining completions all used LLM harness architectures with varying levels of tool integration. Nevertheless, LLMs played a central role in the winning approach: the LLM provided the prior (task structure and initial behavior) while RL provided refinement (faster execution and strategy discovery). This pattern of using LLMs for high-level reasoning and RL for low-level optimization suggests a fruitful direction for future work; we discuss the scalability of LLM--RL hybrid approaches in Section~\ref{sec:future_work}.

The PokeAgent Simple baseline, by contrast was the weakest of the LLM harness approaches. Several observations from its performance spurred subsequent development. First, the release of Gemini 3 models led to a notable increase in performance, particularly around visual processing. Actions appeared more visually grounded, with less overreliance on text state~\citep{google2025gemini3}. Tasks requiring precision tile-level movement, orientation, and interaction, such as setting the clock in the player's bedroom, became much more consistent. This confirmed that base model capability, especially in its ability to natively reason across multiple input modalities, was a first-order driver of performance. Second, the four-module architecture was poor on both step and time efficiency, requiring 3 VLM calls per action, prompting us to increase the efficiency of our action-generation pipeline. Third, the agent lacked coherence across extended play: it would frequently become confused about game progress and infer its next action solely from the current screen and its world knowledge, making slow but eventual progress through the tutorial's largely linear structure. A common failure mode followed the completion of the rival battle in Route 103 (see Figure ~\ref{fig:ch3_speedrun_route} for reference). Upon returning back to Oldale Town, the agent would often attempt to proceed westward to Route 102 without first traveling further south to visit Professor Birch in Little Root Town. Decreasing the prevalence of this failure mode would necessitate a more powerful abstraction for planning.



\subsubsection{Intermediate Harness Improvements}\label{sec:benchmark_results}

Following the competition, the \pokeagent{} harness underwent its first major revision, introducing several architectural changes that substantially improved consistency and speed in the tutorial setting:

\paragraph{Objective System.} The most impactful addition was an explicit objective system that separated planning from action execution. Rather than relying on the underlying model to implicitly maintain a sense of progress through a combination of its world knowledge and action history, the new harness explicitly maintained a sequence of objectives that could be generated dynamically by the agent or pre-defined from a walkthrough source~\citep{bulbapedia2025walkthrough}. To support this new system for planning, we additionally exposed research tools to the agent for obtaining progress and walkthrough summaries. By isolating planning into a distinct phase, the agent could develop structured plans and execute them systematically, reducing the frequency of CRF modes. At this stage, the objective system was an append-only log whose successful employment required accurate completion detection. In the absence of tools enabling replanning, the premature marking of, or failure to denote critical objectives as completed could still lead to unrecoverable states.

\paragraph{Knowledge Base.} A persistent knowledge base abstraction allowed the agent to store discoveries (locations visited, items found, NPCs locations, storyline insights) in a flat list with importance-weighted retrieval. This gave the agent a mechanism for localizing its progress within the game narrative when generating new plans, improving consistency in execution.

\paragraph{A* Pathfinding.} Our initial observations of \pokeagent{} Simple's elucidated the difficulty that even frontier VLMs have in successfully traversing EAEs. Recent evaluations confirm that LLMs struggle with coordinate-based navigation and spatial reasoning tasks, with performance degrading sharply as environment complexity increases~\citep{yamada2024evaluating}, and that agentic LLMs and VLMs exhibit particularly poor navigation performance in game environments~\citep{paglieri2024balrog}.  The integration of A* pathfinding abstracted away most of reasoning burden into a programtic tool. The agent could now request navigation to specific coordinates or warp points, with the pathfinder computing optimal routes and introducing optional randomness that respected tile traversability, collision boundaries, and ledges. This reduced navigation failures dramatically, though the addition of a programmatic tool with such a strong coupling to game state decreased the generality of the harness by imposing an additional engineering burden. Reducing this dependence on hand-crafted pathfinding---for example, through learned spatial representations or SLAM-like map construction---is a direction we discuss in Section~\ref{sec:future_work}.


\subsubsection{Open Benchmark Results}\label{sec:open_benchmark_results}

\begin{figure}[h]
\centering
\includegraphics[width=\linewidth]{ch4/milestone_confidence.pdf}
\caption{\textbf{Speedrunning Harness vs.\ Model Comparisons (Tutorial Setting).} Cumulative wall-clock time, actions, tokens, and cost at each milestone for five frontier models on the \pokeagent{} harness (mean $\pm$ min/max range across runs). Gemini 3 Flash completes the route fastest ($\sim$2:24 mean) but requires more actions than Gemini 3 Pro. Claude Sonnet 4.5 completes all milestones but with the highest variance and 3--4$\times$ the cost of Gemini variants. GPT-5.2 falls between the two families.}
\label{fig:ch4_harness_vs_model}
\end{figure}


\paragraph{Model Comparisons.} A consistent finding across models is that Gemini outperforms in visual reasoning, despite Claude and GPT models being generally regarded as stronger on text-only benchmarks~\citep{chiang2024chatbot, google2025gemini3}. Gemini models exhibited more reliable visual grounding: they would almost always select Mudkip as their starter, with reasoning traces explicitly noting its type advantage against the Rock-type first gym. Claude and GPT models were less consistent in this regard, affecting downstream battle performance and completion speed for the tutorial segment. While not directly relevant to our harness-level analysis---which does not involve training or weight updates---this observation nonetheless suggests that domain-specific world knowledge is a meaningful differentiator for performance even among frontier models.

\paragraph{Harness Comparisons.} The retrospective also evaluated three proprietary CLI agents harnesses: Claude Code, Codex CLI, and Gemini CLI. We provided these harnesses with access to three MCP tools for retrieving game state, pressing buttons, and pathfinding. These coding-focused agents bring powerful native abstractions: context windows on the order of 200,000+ tokens before compaction, strong Pok\'{e}mon world knowledge, and proprietary problem-solving scaffolding that enables hypothesis exploration and online research. Several qualitative observations stood out. Their native planning systems, while effective for software engineering tasks, operated at a level of granularity that was often insufficiently specific for RPG navigation---plans like ``go to Rustboro City'' lack the operational detail needed for reliable execution. Their proprietary scaffolding allowed them to explore hypotheses directly in the embodied setting and use online research tools to redirect when stuck, but at the cost of significantly higher token burn. These results underscored that while powerful general-purpose harnesses can make meaningful progress out of the box---especially when equipped with external tools like pathfinding---the non-linear, long-horizon exploration characteristic of RPG play benefits from domain-tailored abstractions. We return to this observation in Section~\ref{sec:future_work}.

% ---------------------------------------------------------------------------
\subsection{Towards an Expert Harness}\label{sec:towards_subagents}

The results and subsequent improvements made to our harness over the prior two implementations made clear a set of architectural gaps in the harness that could not be addressed by incremental improvements. Two external trends reinforced the direction of development.

\paragraph{Model Capability Improvements.} The transition from Gemini 2.5 to Gemini 3 models brought substantial improvements in multimodal reasoning~\citep{google2025gemini3}. Actions became more visually grounded, with the agent demonstrating finer control over movement and more reliable interpretation of in-game text and navigation of menus. Tasks that had been prohibitively difficult for Gemini 2.5, such as accurately setting the bedroom clock or reading small menu text, became reliable with the Gemini 3 family of models. We likewise observed improvements to the consistency of tool-calling via decreases in the incidences of  passing invalid arguments or erroneously formatted parameters. For instance, Gemini 2.5 models would sometimes call \texttt{PRESS-BUTTONS['QUICK ATTACK']} or \texttt{PRESS-BUTTONS['CLOCK']} rather than using \texttt{PRESS-BUTTONS['A']} when attempting to select a move or interact with a particular item. This progress suggested that a more capable VLM would improve both the agent's intrinsic reasoning and its reliability in using the harness's tool interface.


\paragraph{Subagent Orchestration.} In parallel, the rapid maturation of agent frameworks in the software engineering domain through the likes of Claude Code, Cursor, Codex CLI, and Gemini CLI, demonstrated the power of subagent architectures for managing context in long-horizon tasks~\citep{cursor2026scalingagents}. These systems introduced planning tools, persistent memory, and eventually subagent delegation as first-class primitives. Their success in coding tasks provided a compelling template for embodied agent design, as demonstrated by their out-of-the-box generalization to our speedrunning benchmark. The advent of subagent orchestration, a paradigm wherein a primary orchestrator agent delegates atomized tasks for subagents (equipped with their own context window) to complete, was of particular interest due to its effectiveness in preventing \textit{context rot}\footnote{Context rot refers to the progressive degradation of an agent's reasoning coherence as its context window accumulates irrelevant or contradictory information over many turns. This phenomenon is documented in the context of recursive language models~\citep{zhang2025rlm} and observed directly in long-horizon Pok\'{e}mon play~\citep{gemini2p5report}.}. As we observe, the loss in coherence as a result of context rot is a relevant failure mode within the \pokemon{} Emerald setting where the agent may forget subtle details as it attempts to complete a difficult task over many turns or when switching between overworld, battle, or menu states. Improving the handling of battles was a notable area of improvement as a common failure mode involved the agent mistaking a trainer battle for a wild encounter and wasting resources attempting to catch another trainer's \pokemon{}.

\paragraph{Self-critique.} Emerging work on self-adaptive agent systems demonstrated the value of self-critique abstractions~\citep{madaan2023selfrefine, shinn2023reflexion}. Without such mechanisms, agents cannot meaningfully adjust their policy when stuck, increasing the probability that a failure mode cascades into CRF and requiring manual intervention via system prompt edits (the full prompt is reproduced in Appendix~\ref{sec:appendix_prompts}). This manual approach does not scale to long-horizon experimentation where it is unreasonable to analyze the agent's reasoning traces over multiple hours of execution.

Two recurring failure modes illustrate the need for self-critique. First, the agent was sensitive to the naming and granularity of objectives: an objective like ``route\_104\_south'' would cause the agent to navigate to the southernmost tile of Route~104 rather than recognizing the label as a distinct map section preceding Petalburg Woods. Without a mechanism to detect and correct this misinterpretation, the error would cascade into downstream objective failures. Second, maps with many warp points overwhelmed the agent's ability to select the correct destination, producing loops of navigation to incorrect warps that the agent could not self-diagnose. Both failure modes share a common structure: the agent lacks the ability to step back, evaluate whether its current approach is working, and adjust. As demonstrated by Self-Refine~\citep{madaan2023selfrefine} and Reflexion~\citep{shinn2023reflexion}, equipping agents with mechanisms for verbal self-evaluation can break these failure loops before they become irrecoverable.

These observations motivated the refactoring of the intermediate harness design along several specific abstractions as we looked to improve performance in long-horizon settings. Our expert harness ($\mathcal{H}_{\text{expert}}$) supports (1) subagent delegation for improved context management and defines several custom subagents along with general subagent primitives for orchestrating agent to define at runtime, restructures the Knowledge Base and action-history into (2) persistent hierarchical short and long term memory, introduces (3) a skill library, (4) sandboxed code execution to allow the agent to prototype and deploy their own custom tools, and (5) refactors the objective system to enable the dynamic replanning.

% ---------------------------------------------------------------------------
\subsection{System Design: Expert Harness}\label{sec:system_design}

The expert harness $\mathcal{H}_{\text{expert}}$ is a multi-agent orchestration system and represents a distinct qualitative departure from its preceding architectures: rather than a monolithic observe-plan-act loop, $\mathcal{H}_{\text{expert}}$ decomposes the agent into a central orchestrator that is equipped with a subset ($\mathcal{T_O}$) of the available tools including tools for creating and calling subagents, set of specialized subagents ($\mathcal{S_O}$), persistent stores (subagents, skills, memories), and a sandboxed execution environment. This system is engineered specifically for the long-horizon setting described in Section~\ref{sec:long_horizon}. As we will see, it accepts a modest penalty in per-step efficiency because of the overhead introduced by frequent context switching between orchestrator and subagents, in exchange for more stable performance and reduced incidences of CRF modes, over extended play across longer horizons.

We also formalize three scaffold configurations, at varying levels of sophistication, relevant to the ensuing chapter: \texttt{"\pokeagent{}"} ($\mathcal{H}_{\text{expert}}$) provides the full tool suite including built-in subagents, A* pathfinding, walkthrough access, and wiki lookup; \texttt{"Simple"} ($\mathcal{H}_{\min}$) strips away all built-in subagents, navigation tools, and research tools, leaving only raw button presses and the persistent stores; \texttt{"AutoEvolve"} ($\mathcal{H}_{\text{auto}}$) extends $\mathcal{H}_{\min}$ with a \texttt{HarnessEvolver} that mutates skills, subagents, and system prompts over time. The latter two configurations are central to the experimental design in Chapter~\ref{sec:autoevolve}.

\begin{figure}[t]
\centering
\includegraphics[width=\linewidth]{ch4/pokeagent_architecture.pdf}
\caption{\textbf{$\mathcal{H}_{\text{expert}}$ Multi-Agent Architecture.} The orchestrator coordinates subagents and MCP tools to play the game. Data flows left-to-right from the GBA emulator (game frames, parsed state) through the orchestrator (context management, subagent dispatch, stuck detection), which invokes either MCP tools (pathfinding, button inputs, knowledge retrieval) for direct game control or specialized subagents (battler, reflection, verification, gym puzzle, planner) for complex reasoning. A feedback loop returns compacted subagent outputs and tool results to inform subsequent decisions.}

\tersoo{TODO: Refine this draft with the full expert harness details (trajectory window, memory/skill/subagent stores, sandbox).}
\label{fig:ch4_expert_architecture}
\end{figure}

\subsubsection{Architecture Overview}\label{sec:architecture_overview}

At each step, the orchestrator fetches the current game state, assembles a structured prompt incorporating the short-term memory window, persistent store overviews, current objectives, and the visual frame, then queries the VLM and executes any resulting tool calls (Figure~\ref{fig:ch4_expert_architecture}). Every step must conclude with either a \texttt{press\_buttons} or \texttt{navigate\_to} call, preventing the model from deliberating indefinitely without advancing the game. When the orchestrator delegates to a subagent (e.g., the battler or planner), control transfers to a dedicated handler that runs its own multi-turn loop with a restricted tool set and fresh game state each turn. Upon termination, only a compacted summary returns to the orchestrator's short-term memory---not every inner turn. This context isolation prevents the orchestrator's context window from accumulating hundreds of battle or planning turns, enabling it to maintain coherence over thousands of high-level decisions. The following subsections detail each abstraction; a complete tool reference is provided in Appendix~\ref{sec:appendix_tools}.

\subsubsection{Trajectory Window}\label{sec:trajectory_window}

The trajectory window is a persistent global session log that records every agent step as a structured entry containing both the meta data associated with an agent step (step number, action taken, reasoning trace, outcome, location, and player coordinates) as well as full prompt passed to the VLM at that step---including the constituent structured text state. While it is also possible to include the image frame at these steps, the trajectory window that we employ does not save this information. This log takes inspiration from the treatmeant of global context by the RLM paradigm~\citep{qu2024recursive}. Rather than viewing the agent's history history as a flat, append-only context sequence, the trajectory is a structured, programmatically accessible workspace that the orchestrator process with the aid of subagents to query, analyze, and summarize any segment of execution history. We expose a tool (\texttt{process\_trajectory\_history}; see Table~\ref{tab:tool_defs} in Appendix~\ref{sec:appendix_tool_defs}) that allows the orchestrator to run ad-hoc analysis over arbitrary step ranges, enabling the agent to reason about patterns in its own behavior (e.g, verifying progress, reflecting on performance to unblock itself, etc).

% [Content moved to Chapter 5 -- trajectory window as bridge to AutoEvolve]


\subsubsection{Memory}\label{sec:memory}

The memory system implements a dual-layer architecture loosely inspired by the distinction between working memory and long-term memory drawn in cognitive neuroscience ~\citep{baddeley2000episodic, tulving1972episodic}. Reading from the short-term buffer---akin to human working memory--- is fast and automatic, while querying the long-term store requires an explicit tool call at higher latency in a way that mirrors the effortful retrieval from long-term memory. 

\paragraph{Short-term memory.} Both the orchestrator and any looping subagent internally maintain a buffer of the previous 20 tool calls---their parameters, success, and embedded agent reasoning. This buffer, formatted and injected into every prompt as the \textsc{Short-Term Memory} section, provides immediate context about recent actions and their outcomes. Additionally, compact overviews, organized as file-trees, of all three persistent stores (memory, skills, subagents) are injected every step, giving the orchestrator a birds-eye view of its accumulated knowledge without consuming excessive context.

\paragraph{Long-term memory.} The persistent memory store contrasts short-term memory by virtue of its tradeoff between increased latency for a virtually unbounded amount of data storage. This store supports CRUD operations with hierarchical path organization. Each \texttt{MemoryEntry} contains an identifier, a hierarchical path (e.g., \texttt{pokemon/team}, \texttt{events/gym\_battles}), a title, free-text content, the location and coordinates where it was created, importance (1--5), and tags. The tree overview renders entries as an indented hierarchy with \texttt{[id] title} leaves, providing a compact index that the VLM can scan to decide which entries to read into its short-term memory, update, add, or delete. Once a long-term memory is encoded within the short term memory as a tool-call result, it propagates there for a number of turns proportional the size of the short-term memory buffer. This asymmetry incentivizes the agent to maintain a compact, well-organized representation of its global knowledge, enabling the efficient retrieval of detailed knowledge.

\subsubsection{Skills}\label{sec:skills}

The global skill library extends the same process of storage and retrieval to behavioral and executable strategies that are learned at runtime and accessible to both the orchestrator and any subagents. Each skill contains a name, description, hierarchical path, empirical effectiveness rating (low/medium/high), and optionally a field containing executable Python code. The tree overview annotates executable code skills with \texttt{(run with: run\_skill)}, signaling to the VLM that these skills can be invoked directly.

Behavioral skills are textual: they encode strategies (e.g., ``when facing Rock-type gym leaders, lead with a Water or Grass type and use super-effective moves'') that an agent may optionally consult during reasoning. Executable skills contain Python code that runs in the sandbox (Section~\ref{sec:sandbox}) with access to game tools, enabling the agent to automate repetitive patterns---such as a training its team, navigation, or item acquistion---without replanning from scratch.

Agents manages skills in a similar fashion to memories as our harness exposes \texttt{process\_skill} (CRUD) and \texttt{run\_skill} (execution). 

% [Content moved to Chapter 5 -- skill library population in H_expert vs H_auto]

\subsubsection{Subagents}\label{sec:subagents}

The subagent system represents the most substial architectural divergence from prior iterations of our agent harness. It distinguishes between two types of subagents: \emph{1-step} subagents that perform a single VLM pass over the trajectory window without tool access, and \emph{delegated-loop} subagents that are callable from the orchestrator and run as autonomous multi-turn loops with their own restricted tool sets.

\paragraph{Built-in 1-step subagents.} $\mathcal{H}_{\text{expert}}$ exposes four pre-configured 1-step subagents that operate on the trajectory window and provide structured guidance (prompt text in Appendix~\ref{sec:appendix_prompts}):
\begin{itemize}[nosep]
    \item \texttt{reflect}: an external critique agent prompted to diagnoses stuck or looping states by analyzing the trajectory, returning an assessment with specific recommendations for strategy changes.
    \item \texttt{verify}: independently checks whether the current objective is truly complete, returning structured evidence for and against completion with a confidence score. This external oversight helps prevent the orchestrator from prematurely advancing when the agent incorrectly believes a task is finished.
    \item \texttt{summarize}: produces an unbiased trajectory summary, used for context compaction during context switches between the orchestrator and predefined subagents to preserve essential information while discarding redundant detail. Agents can optionally call this summarization tool ad-hoc.
    \item \texttt{gym\_puzzle}: provides domain-specific guidance for solving gym puzzles (e.g., Mauville's electric barriers) combined with current-state analysis.
\end{itemize}

\paragraph{Built-in delegated-loop subagents.} Two subagents run autonomous multi-turn loops:
\begin{itemize}[nosep]
    \item \texttt{battler}: handles all combat decisions---move selection, switching, item usage---for up to 200 turns with a restricted tool set (Table~\ref{tab:tool_access}). The orchestrator delegates control when a battle is detected and receives a compacted summary upon completion.
    \item \texttt{planner}: manages objective creation and revision for up to 25 turns, with access to research tools (Table~\ref{tab:tool_access}). The planner sees the full objective sequence across all categories (story, battling, dynamic) and can create, modify, or delete objectives.
\end{itemize}

\paragraph{Custom subagent primitive.} Beyond the built-in subagents, the orchestrator can create and launch arbitrary subagents at runtime via two mechanisms: (1)~registering a new subagent in the persistent subagent registry with a specified handler type, maximum turns, available tools, system instructions, and return condition, then launching it by ID; or (2)~launching an inline subagent with an ephemeral configuration. Custom subagents run through the same executor infrastructure as built-ins, with a \texttt{return\_to\_orchestrator} signal to handle termination. The full tool access matrix for all subagent types is provided in Table~\ref{tab:tool_access}.

In both custom and predefined subagents, both nested calls in which subagents are called recursively and parallel subagent execution are unsupported.

% [Content moved to Chapter 5 -- custom subagent primitive and AutoEvolve expressivity]

\subsubsection{Sandbox and Code Execution}\label{sec:sandbox}

Two tools support sandboxed Python execution:

\paragraph{\texttt{run\_skill}} loads a skill's \texttt{code} field from the skill store and executes it in a sandbox with restricted built-in functions (no \texttt{open}, \texttt{os}, \texttt{sys}, \texttt{subprocess}) and access to standard-library modules (\texttt{random}, \texttt{collections}, \texttt{math}, \texttt{json}, \texttt{re}, \texttt{heapq}, \texttt{itertools}, \texttt{functools}, \texttt{numpy}). Game interaction is mediated through a \texttt{tools} dictionary providing callable wrappers for \texttt{press\_buttons}, \texttt{get\_game\_state}, \texttt{get\_map\_data}, \texttt{complete\_direct\_objective}, and \texttt{process\_memory}. The skill's return value is passed back to the VLM as the tool result.

\paragraph{\texttt{run\_code}} provides the same sandbox but with read-only tools only (\texttt{get\_game\_state}, \texttt{get\_map\_data})---no \texttt{press\_buttons}. It captures \texttt{print()} output to a string buffer, returning both the printed output and any result value. This tool is designed for prototyping and debugging: the agent can test code logic before committing it to the skill store as a reusable executable skill.

The sandbox is process-local (not containerized), with security enforced through restricted built-ins. The \texttt{\_\_import\_\_} function remains accessible, allowing the agent to import additional standard-library modules as needed, but game-action tools are explicitly gated to prevent uncontrolled execution.

\subsubsection{Objective System}\label{sec:objective_system}

The intermediate harness's append-only objective log (Section~\ref{sec:benchmark_results}) was replaced with a CRUD-based objective system. Objectives are organized into three categories---\emph{story} (narrative progression milestones), \emph{battling} (team composition and training targets), and \emph{dynamics} (agent-generated situational goals)---each maintained as an ordered sequence. The planner subagent accesses the full objective state via \texttt{get\_full\_objective\_sequence} and modifies it through \texttt{replan\_objectives}, which supports index-based additions, completions, reorderings, and deletions within a single category per call. This replanning capability addresses the failure mode identified in the intermediate harness: prematurely completed or incorrectly ordered objectives can now be corrected without manual intervention, making the system resilient to the cascading errors that previously led to CRF.

\vspace{1em}
Taken together, these six abstractionsx define a harness that is substantially more capable than the competition-era agent. The cost is the number and interdependence of these components: each abstraction interacts with several others (e.g., skills execute in the sandbox with access to memory; subagents query the trajectory window and modify objectives), and the orchestrator must coordinate all of them coherently across thousands of steps. Whether this engineering effort can be partially or fully automated is the central question of Chapter~\ref{sec:autoevolve}.
